\section{A pattern-preserving affine corner trade}
\label{sec:corner-trade}

Latin squares with a unique intercalate are an established subject;
see Wanless~\cite{Wanless2001Unique}. The result here concerns the
three-view FFF pattern of a particular affine trade. Neither the
existence of unique-intercalate squares nor priority for this formula
is claimed to be new.

For an odd prime $p$ and $a\in\mathbb F_p\setminus\{0,-1\}$, define
$P_a$ on $\mathbb F_p\cup\{\infty\}$ by
\[
 P_a(r,c)=\begin{cases}\infty&c=ar,\\r+c&c\ne ar\end{cases}
 \quad(r,c\in\mathbb F_p),
\]
with $P_a(r,\infty)=(a+1)r$, $P_a(\infty,c)=(a+1)c/a$, and
$P_a(\infty,\infty)=\infty$. This is the standard one-point
prolongation of addition along its affine transversal. Let $P'_a$
swap the two entries in each of rows $0,\infty$ at columns $0,\infty$.
Those four cells form an intercalate.

\begin{theorem}[Corner trade]
\label{thm:corner-trade}
Both tables are Latin and $\pat(P'_a)=\pat(P_a)$ componentwise.
If $p\geq5$ and $a\notin\{1,-2,-1/2\}$, then
$I(P_a)=p$ and $I(P'_a)=1$.
\end{theorem}

\begin{proof}
The prolongation is Latin, and an intercalate swap preserves this property.
Put $b=(a+1)/a$, $d=\operatorname{ord}_p(b)$ and let $a_0$ be the
integer representative of $a$ in $\{1,\ldots,p-2\}$. The autotopisms
\[
 (r,c,s)\longmapsto(r+t,c+at,s+(a+1)t)
\]
fix infinity. Together with the scaling autotopisms
$(r,c,s)\mapsto(ur,uc,us)$ for $u\ne0$, they make all distinct
finite-finite row pairs have the same cycle type. For rows $(0,t)$
with $t\ne0$, scaling columns by
$t^{-1}$ gives the relative map
\[
 0\mapsto a,\quad \infty\mapsto-1,\quad a+1\mapsto\infty,
 \qquad x\mapsto x-1\quad(x\ne0,a+1).
\]
Its cycles containing $0$ and infinity are separate, of lengths
$a_0+1$ and $p-a_0$. For rows $(\infty,t)$, the ordinary finite
map is $x\mapsto bx-t$. Prolongation replaces its fixed point $at$
by $(at,\infty)$; the remaining cycles have length $d$. If $t\ne0$,
the point $0$ belongs to one of those $d$-cycles.

Writing $\tau=(0,\infty)$ on columns, the trade precomposes rows
$0,\infty$ with $\tau$. The pair $(0,\infty)$ is conjugated and
keeps its type. For $(0,t)$ the transposition joins the two cycles
into one $(p+1)$-cycle. For $(\infty,t)$ it joins the two-cycle
containing infinity to a separate $d$-cycle, giving lengths
$d+2,d,\ldots,d$. Pairs of distinct nonzero finite rows are unchanged.
Thus row-F holds after the trade exactly when $a_0+1,p-a_0,d$ are
all even, which is also its condition before the trade. Necessity
uses an unchanged nonzero finite pair and the pair $(0,\infty)$;
these exist also for $p=3$.

Transposition gives exactly $P_a^T=P_{a^{-1}}$, preserving the
corner trade. In the symbol view, take symbols as lines and columns
as positions, and relabel finite columns by $c'=-c$. Solving the
table equation gives precisely the same construction with parameter
$-a/(a+1)$, again taking the corner to the corner. Applying the row
calculation to these two parameters proves column and symbol equality.

Each two-cycle in a row-pair permutation is one intercalate. In $P_a$
the $p$ pairs involving infinity each contribute their distinguished
two-cycle. Additional two-cycles occur exactly when $d=2$, $a_0+1=2$
or $p-a_0=2$, that is, $a=-1/2,1,-2$. Outside these exceptions
the count is $p$. After the trade, the complete pair list above
leaves only the distinguished two-cycle of $(0,\infty)$, giving one.
\end{proof}

\begin{corollary}[One intercalate and no standard two-point contraction]
\label{cor:corner-noncontraction}
For every prime $p\geq10^{12}$ there exists a non-group-isotopic FFF
square of order $p+1$ with exactly one intercalate. It is not isotopic
or parastrophic to a standard two-transversal prolongation of an
order-$(p-1)$ square with all retained cells unchanged.
\end{corollary}

\begin{proof}
The isotopy $x=(a+1)r$, $y=(a+1)c/a$, with symbols unchanged,
takes $P_a$ to the loop of Section~\ref{sec:affine-prime} with parameter
$\alpha=1/(a+1)$. The inverse correspondence
$a=(1-\alpha)/\alpha$ is bijective. Theorem~\ref{thm:affine-large-prime}
gives more than $p/1000>3$ FFF parameters; at most three are exceptional
for Theorem~\ref{thm:corner-trade}. This proves existence.
A group table of even order $n$ with $t$ nonidentity involutions has
$tn^2/4$ intercalates: each involution gives $n/2$ unordered row pairs
with $n/2$ two-cycles. Here $t\geq1$ and $n\geq4$, so this is greater
than one. Intercalate count is isotopy invariant.

The unique intercalate of a generic $P'_a$ has rows, columns and symbols
$\{0,\infty\}$. Delete them. For each $r\ne0$, the retained position
$(r,-r)$ is a hole because its symbol was zero. Its possible values
must lie in
\[
 \{r,(a+1)r\}\ \cap\ \{-r,-(a+1)r/a\}.
\]
The four possible equalities respectively force characteristic two,
$a=-1/2$, $a=-2$ or $a=-1$, all excluded. Thus the intersection is
empty and even one such hole forbids a Latin completion with retained
entries fixed. A standard two-transversal prolongation has an order-two
corner whose reverse operation is just such a completion. Uniqueness
of the intercalate and invariance under all three-sort relabellings and
coordinate permutations prove the final assertion.
\end{proof}

This corollary inherits the published character-sum estimate used in
Theorem~\ref{thm:affine-large-prime}; it is not a physical computation
at enormous orders. Independently scanned small examples are
$(p,a)=(13,3)$ and $(19,7)$, of orders $14$ and $20$. The finite
checker tests pattern preservation at all 253 admissible parameters
over primes through $43$ and the intercalate and empty-domain formulas
at all 216 generic parameters. Every one of the 15 affine cases at
$p=17$ remains non-FFF. No arbitrary order-$18$ claim follows.

\begin{computationaltheorem}[Expanded explicit class lower bounds]
\label{cthm:fff20-1703}
There are at least $1703$ FFF main classes at order $20$, at least
$T_d(1703)$ at order $20\cdot2^d$, and in particular at least
$1450956$ at order $40$.
\end{computationaltheorem}

\begin{proof}
The frozen enlargement of the Steiner input tests every trade mask
for each of its 22 all-two-cycle row pairs. Its historical enumeration
records 22528 masks, 19480 FFF labelled mask instances and 1700 distinct
exact aggregate spectra. Two quadratic controls add two further spectra.
For the present lower bound it suffices to reconstruct one specified
table per spectrum: the fresh companion checker reconstructs all 1702,
validates Latinness and all three views by two methods, and compares
the spectra as exact tuples, not merely hashes. Each has at least 19
intercalates; the checked $P'_7$ at $p=19$ has one and supplies another
class.

For all 1703 selected tables, fresh binary elimination produces a
$58$-square minor, and a separate explicit matrix elimination checks
its nonsingularity and binding to the table. Lemma~\ref{lem:binary-rank-certificate}
establishes binary quotient freeness without C157. Apply
Theorem~\ref{thm:affine-fibre-orbits} and
Equation~\eqref{eq:burnside-fibres}; at $d=1$ the count is
$1703\cdot1704/2=1450956$.
\end{proof}

The fresh representative audit proves these lower bounds without
repeating the complete 22528-mask search. It does not prove an exact
number of all FFF classes, completeness of the spectrum invariant,
or literature priority. The older 514-input bounds remain valid but
are no longer the strongest bound in this edition.
